\part{Introduction to Security}

Every attack in this book eventually reduces to the same trick: get a system, or the person
running it, to do something it should not. Unit 1 starts with the person, because the person is
almost always the softer target. A firewall does not get tired at 4:45 on a Friday, does not want
to seem helpful to a caller who claims to be from IT, and does not have a mortgage payment that
makes a ``verify your account'' email feel urgent instead of suspicious. Chapter~\ref{ch:1.1}
opens with social engineering for exactly this reason: before any technical control matters, the
attacker has to get past a human, and most of the time that human has never been taught what the
attempt looks like.

From there the unit moves outward, from the person to the systems the person touches every day.
Chapter~\ref{ch:1.2} asks how a login page can lie to you, and what a browser is and is not able
to guarantee about the site on the other end. Chapter~\ref{ch:1.3} moves off the device entirely
and onto the network in between, where a coffee-shop Wi-Fi connection turns every unencrypted
packet into something a stranger three tables away can read. The last two chapters bring in the
newest variable in this whole equation: artificial intelligence. Chapter~\ref{ch:1.4} covers how
AI lowers the cost of running a convincing social-engineering or phishing campaign at scale, and
Chapter~\ref{ch:1.5} closes the unit by turning the same technology around, using it to detect the
attacks it just got easier to launch.

None of this requires deep technical background yet, which is deliberate. The goal of Unit 1 is
to build a working vocabulary and a healthy suspicion of anything that asks for trust before it
has earned it, since every later unit assumes that instinct is already in place.
